\chapter{Tests et validation}

Pour garantir que le site web est exempt de bugs en production et qu'il offre une expérience utilisateur irréprochable, j'ai mis en place une stratégie de validation divisée en trois axes : Tests Fonctionnels, Tests de Performance, et Tests de Sécurité.

\section{La Pyramide des Tests}
Pour garantir une stabilité logicielle professionnelle, j'ai implémenté le concept de la "Pyramide des Tests", qui consiste à avoir une large base de tests rapides, complétée par des tests plus lourds simulant le comportement humain.

\begin{itemize}
    \item \textbf{Tests Unitaires (Jest) - La base de la pyramide :} J'ai écrit des dizaines de scripts automatisés pour tester isolément chaque fonction mathématique et métier. Par exemple, la fonction de calcul \texttt{calculateTotalTTC(prixHT, jours, codePromo)} est testée avec diverses entrées extrêmes (valeurs négatives, décimales infinies) pour s'assurer que le système comptable du \texttt{service-erp} ne génère jamais de fausses factures.
    \item \textbf{Tests d'Intégration (Postman / Supertest) - Le milieu de la pyramide :} Il s'agit de vérifier que le \texttt{service-location} parvient bien à communiquer avec la base MongoDB. J'ai simulé des requêtes HTTP POST pour m m'assurer que l'API renvoie bien le code d'erreur HTTP 401 (Unauthorized) si un utilisateur tente de réserver un véhicule avec un Token JWT expiré.
    \item \textbf{Tests End-to-End (Cypress) - Le sommet de la pyramide :} Cypress se comporte comme un "robot" simulant un véritable utilisateur humain. Le script lance un navigateur Chrome fantôme, clique sur "Se connecter", tape des identifiants au clavier, choisit une voiture, clique sur "Payer", et vérifie visuellement que le message de succès apparaît à l'écran. Cela garantit que toute la chaîne (du clic HTML jusqu'à l'écriture en BDD) fonctionne de bout en bout.
\end{itemize}

\section{Tests de Charge et de Performance}
Outre les fonctionnalités, une plateforme web doit pouvoir résister à de forts pics d'affluence (par exemple, à l'approche de la saison estivale).
\begin{itemize}
    \item \textbf{Stress Test (JMeter / K6) :} J'ai soumis l'API Node.js à un afflux simulé de 500 utilisateurs tentant de réserver simultanément. Grâce à l'architecture asynchrone de Node.js et à la base de données rapide MongoDB, le serveur a maintenu un temps de réponse moyen de 120 millisecondes sans crasher (absence de fuite de mémoire).
    \item \textbf{Validation des Performances Front-End (Google Lighthouse) :} J'ai audité le portail React (\texttt{enterprise-portal}) pour optimiser les \textit{Core Web Vitals}. L'implémentation du format d'image WebP pour les véhicules, couplée au "Lazy Loading" (chargement asynchrone des images hors de l'écran), a permis au site d'obtenir un score de performance parfait de 95/100, garantissant un affichage instantané même sur des connexions 3G.
\end{itemize}

\section{Tests de Sécurité (DevSecOps)}
Dans le cadre de l'intégration continue (CI/CD) sur GitHub Actions, j'ai intégré un scan automatisé avec \textbf{SonarQube}. Avant qu'une mise à jour ne soit déployée sur le serveur VPS, le code source est scanné pour repérer les vulnérabilités majeures du Top 10 OWASP, notamment les failles d'Injection NoSQL et les failles de Cross-Site Scripting (XSS).

\section{Cas de Tests (Tableau Récapitulatif)}

\begin{table}[H]
    \centering
    \renewcommand{\arraystretch}{1.5}
    \begin{tabular}{|p{4cm}|p{4cm}|p{4cm}|p{2cm}|}
        \hline
        \cellcolor{gray!20}\textbf{Scénario de test} & \cellcolor{gray!20}\textbf{Résultat attendu} & \cellcolor{gray!20}\textbf{Résultat obtenu} & \cellcolor{gray!20}\textbf{Statut} \\
        \hline
        Tentative de connexion avec un mauvais mot de passe & Message d'erreur "Identifiants incorrects" affiché & Message affiché, token non délivré & \textcolor{green}{SUCCÈS} \\
        \hline
        Réservation d'un véhicule déjà loué & Rejet immédiat avec erreur 409 Conflict & Erreur de conflit gérée par le verrou optimiste & \textcolor{green}{SUCCÈS} \\
        \hline
        Génération d'une facture par un simple "Commercial" & Accès refusé (Forbidden 403) & Middleware RBAC bloque la route & \textcolor{green}{SUCCÈS} \\
        \hline
    \end{tabular}
    \caption{Cas de tests de l'application}
\end{table}
